%==============================================================================
\chapter{Transportverschijnselen op moleculaire schaal}
\label{ch:transport}
%==============================================================================

\section{Voorbij het thermodynamisch evenwicht}

In de vorige hoofdstukken hebben we de thermodynamica ontwikkeld als een beschrijving van systemen in evenwicht, of systemen die evolueren tussen opeenvolgende evenwichtstoestanden. De thermodynamica voorspelt de begin- en eindtoestand, maar zegt niets over \emph{hoe} het systeem van de ene toestand naar de andere evolueert, noch \emph{hoe snel} de overgang verloopt.

In werkelijkheid verlopen vrijwel alle industri\"ele processen buiten het thermodynamisch evenwicht. In een warmtewisselaar stroomt warmte van een hete naar een koude processtroom met een eindige snelheid. In een reactor diffunderen reactanten naar de katalysator en reactieproducten diffunderen weg. In een destillatiekolom verplaatst materie zich van de ene tray naar de andere. In al deze processen is sprake van \emph{transport}: overdracht van warmte, materie of impuls van een region naar een andere, gedreven door een gradi\"ent in temperatuur, concentratie of snelheid.

Dit hoofdstuk onderzoekt deze transportverschijnselen op moleculaire schaal. We beginnen met de gemeenschappelijke thermodynamische oorsprong (entropieproductie), behandelen de drie fundamentele transportwetten (diffusie, warmtegeleiding, viskeuze stroming), en werken verder naar praktische correlaties zoals Stokes--Einstein. De nadruk ligt op de fysische basis; de formele reactor- en equipmentontwerp-berekeningen worden in een vervolgcursus behandeld.

\section{Entropieproductie in irreversibele processen}

In Hoofdstuk~3 zagen we dat voor een irreversibel proces de totale entropie van het universum toeneemt: $\Delta S_{\text{univ}} > 0$. Voor een lokaal transportproces kan deze entropieproductie worden gekwantificeerd als het product van een \emph{flux} (stroom per oppervlakte-eenheid) en een \emph{kracht} (gradi\"ent in een toestandsvariabele):
\begin{equation}
\sigma_{\text{entropie}} = \sum_i J_i \cdot X_i \geq 0,
\end{equation}
waarin $J_i$ de flux (bv. warmtestroom, massastroom) en $X_i$ de thermodynamische kracht (bv. temperatuurgradi\"ent, concentratiegradi\"ent) is. Het gelijkteken geldt enkel in evenwicht (alle gradi\"enten nul).

Deze formulering verenigt alle transportprocessen onder een gemeenschappelijke noemer: warmtestrooming, diffusie, viskeuze wrijving en elektrische stroom zijn allen uitingen van hetzelfde principe --- het streven naar maximal entropie. Voor de ingenieur is dit een belangrijk inzicht: alle transportprocessen delen dezelfde wiskundige structuur, hetgeen de analyse van complexe systemen aanzienlijk vereenvoudigt.

\section{De drie fundamentele transportwetten}

De flux van een grootheid is in eerste benadering evenredig met de gradi\"ent van de corresponderende toestandsvariabele. Drie equivalente wetten beschrijven dit gedrag.

\subsection{Wet van Fourier (warmtegeleiding)}

De warmtestroom per oppervlakte-eenheid $q$ is evenredig met de temperatuurgradi\"ent:
\begin{equation}
\boxed{\;q = -\lambda\,\frac{\dx{T}}{\dx{x}}\;}
\end{equation}
waarin $\lambda$ de \emph{warmtegeleidingsco\"effici\"ent} is (\si{\watt\per\metre\per\kelvin}). Het minteken garandeert dat warmte van hoog naar laag temperatuur vloeit.

Typische waarden bij \SI{300}{\kelvin}:
\begin{itemize}
\item Koper: $\lambda \approx \SI{400}{\watt\per\metre\per\kelvin}$ (uitstekende geleider).
\item Water: $\lambda \approx \SI{0.6}{\watt\per\metre\per\kelvin}.
\item Lucht: $\lambda \approx \SI{0.026}{\watt\per\metre\per\kelvin} (slechte geleider --- isolatie).
\item Polyurethaanschuim: $\lambda \approx \SI{0.025}{\watt\per\metre\per\kelvin}.
\end{itemize}

\subsection{Wet van Fick (diffusie)}

De massastroom per oppervlakte-eenheid $j$ is evenredig met de concentratiegradi\"ent:
\begin{equation}
\boxed{\;j = -D\,\frac{\dx{c}}{\dx{x}}\;}
\end{equation}
waarin $D$ de \emph{diffusieco\"effici\"ent} is (\si{\metre\squared\per\second}). Typische waarden: gas in gas $D \sim \SI{1e-5}{\metre\squared\per\second}$, vloeistof in vloeistof $D \sim \SI{1e-9}{\metre\squared\per\second}$, vaste stof $D \sim \SI{1e-12}{\metre\squared\per\second}$ of kleiner.

\subsection{Wet van Newton (viskeuze schuifspanning)}

De schuifspanning $\tau$ in een bewegende vloeistof is evenredig met de snelheidsgradi\"ent:
\begin{equation}
\boxed{\;\tau = -\eta\,\frac{\dx{v}}{\dx{x}}\;}
\end{equation}
waarin $\eta$ de \emph{dynamische viscositeit} is (\si{\pascal\second}). Typische waarden bij \SI{20}{\celsius}: water $\eta \approx \SI{1.0e-3}{\pascal\second}; lucht $\eta \approx \SI{1.8e-5}{\pascal\second}; glycerol $\eta \approx \SI{1.0}{\pascal\second}.

\begin{voorbeeld{Warmteverlies door een isolatie}
Een rechthoekige isolatielaag polyurethaanschuim ($\lambda = \SI{0.025}{\watt\per\metre\per\kelvin}$) met dikte $L = \SI{5}{\centi\metre}$ scheidt een binnenmuur ($T_1 = \SI{20}{\celsius}$) van een buitenmuur ($T_2 = \SI{-10}{\celsius}$). De warmtestroom per oppervlakte-eenheid bedraagt:
\[
q = -\lambda\,\frac{T_2 - T_1}{L} = -(0{,}025)\,\frac{(-10) - 20}{0{,}05} = \SI{15}{\watt\per\metre\squared}.
\}
Voor een muur van \SI{10}{\metre\squared} bedraagt het warmteverlies \SI{150}{\watt}. Verdubbeling van de dikte halveert dit verlies --- een directe toepassing van de wet van Fourier.
\end{voorbeeld}

\section{De vergelijkingen van transport}

Wanneer we de fluxwetten combineren met een massabalans of energiebalans, ontstaan differentiaalvergelijkingen die de tijds- en ruimtelijke evolutie van een veld beschrijven.

\subsection{De diffusievergelijking}

Voor diffusie in \'e\'en dimensie geldt:
\begin{equation}
\boxed{\;\frac{\partial c}{\partial t} = D\,\frac{\partial^2 c}{\partial x^2}\;}
\end{equation}
Deze vergelijking beschrijft hoe een concentratieprofiel in de tijd evolueert. Een karakteristieke tijdschaal voor diffusie over afstand $L$ is:
\begin{equation}
\boxed{\;t_{\text{diff}} \approx \frac{L^2}{D}\;}
\end{equation}
Dit is een van de belangrijkste schaalwetten in de fysische chemie.

\begin{voorbeeld{Diffusietijdschaal in cellen}
Een zuurstofmolecule diffundeert in water met $D \approx \SI{2.0e-9}{\metre\squared\per\second}$. De tijd om een afstand $L = \SI{10}{\micro\metre}$ (typische celgrootte) te overbruggen:
\[
t_{\text{diff}} = \frac{(\SI{10e-6}{\metre})^2}{\SI{2.0e-9}{\metre\squared\per\second}} = \frac{\SI{1e-10}{\metre\squared}}{\SI{2.0e-9}{\metre\squared\per\second}} = \SI{0.05}{\second}.
\}
Voor een afstand $L = \SI{1}{\milli\metre} (1000 maal groter):
\[
t_{\text{diff}} = \frac{(\SI{1e-3}{\metre})^2}{\SI{2.0e-9}{\metre\squared\per\second}} = \SI{500}{\second} \approx \SI{8.3}{\minute}.
\}
De kwadratische afhankelijkheid van $L$ is cruciaal: \emph{verdubbeling van de afstand vereist viermaal zoveel tijd}. Dit verklaart waarom biologische cellen klein zijn en bloedvaten een dicht netwerk vormen --- langzame diffusie zou anders de cel van zuurstof beroven.
\end{voorbeeld}

\subsection{De warmtevergelijking}

Analoog aan diffusie geldt voor warmtegeleiding in een medium met constante eigenschappen:
\begin{equation}
\frac{\partial T}{\partial t} = \alpha\,\frac{\partial^2 T}{\partial x^2}, \qquad \alpha = \frac{\lambda}{\rho c_p},
\end{equation}
waarin $\alpha$ de \emph{temperatuurleidingsco\"effici\"ent} is (\si{\metre\squared\per\second}). De karakteristieke tijdschaal is $t_{\text{therm}} \approx L^2/\alpha$. Voor water bedraagt $\alpha \approx \SI{1.4e-7}{\metre\squared\per\second}, hetgeen verklaart waarom warmte zich in water trager verspreidt dan materie.

\section{De oorsprong van transportco\"effici\"enten}

De transportco\"effici\"enten $\lambda$, $D$ en $\eta$ zijn geen empirische constanten maar volgen uit de moleculaire structuur. De kinetische theorie van gassen geeft voor een ideaal gas:
\begin{align}
\eta_{\text{gas}} &\approx \tfrac{1}{3}\,\rho\,\bar{v}\,\ell, \\
\lambda_{\text{gas}} &\approx \tfrac{1}{3}\,\rho\,c_v\,\bar{v}\,\ell, \\
D_{\text{gas}} &\approx \tfrac{1}{3}\,\bar{v}\,\ell,
\end{align}
waarin $\bar{v}$ de gemiddelde moleculaire snelheid en $\ell$ de gemiddelde vrije weg is. Deze vergelijkingen verbinden macroscopische transportco\"effici\"enten met microscopische grootheden --- een directe toepassing van de brug tussen macro en micro die in Hoofdstuk~2 en Hoofdstuk~3 werd gelegd.

\subsection{Gemiddelde vrije weg}

De gemiddelde vrije weg $\ell$ is de gemiddelde afstand die een molecule aflegt tussen twee opeenvolgende botsingen. Voor een ideaal gas:
\begin{equation}
\ell = \frac{1}{\sqrt{2}\,\pi d^2 n},
\end{equation}
waarin $d$ de molecuuldiameter en $n$ de molecuuldichtheid is. Bij atmosferische druk bedraagt $\ell \sim \SI{70}{\nano\metre}$ voor lucht; bij \SI{1}{\milli\bar} (vacu\"um) stijgt $\ell$ tot \SI{70}{\micro\metre}, wat kritisch is voor halfgeleiderproductie en dunnefilmdragers.

\section{De vergelijking van Stokes--Einstein}

Voor diffusie van een bolvormig deeltje (straal $r$) in een vloeistof met viscositeit $\eta$ geldt de \emph{vergelijking van Stokes--Einstein}:
\begin{equation}
\boxed{\;D = \frac{k_B T}{6\pi \eta r}\;}
\end{equation}
Deze vergelijking verbindt diffusie (macroscopisch) met thermische energie $k_B T$ en vloeistofviscositeit (microscopisch). Een direct gevolg: grotere deeltjes diffunderen trager; hogere temperatuur versnelt diffusie; hogere viscositeit vertraagt diffusie.

\begin{voorbeeld{Diffusie van albumine in bloed}
Albumine is een bloedeiwit met straal $r \approx \SI{3.5}{\nano\metre}. In bloedplasma ($\eta \approx \SI{1.5e-3}{\pascal\second}$) bij \SI{310}{\kelvin}:
\[
D = \frac{(\SI{1.38e-23}{\joule\per\kelvin})(\SI{310}{\kelvin})}{6\pi (\SI{1.5e-3}{\pascal\second})(\SI{3.5e-9}{\metre})} \approx \SI{1.4e-10}{\metre\squared\per\second}.
\}
De diffusietijd over \SI{10}{\micro\metre} (afstand tot een haarvat) bedraagt:
\[
t = \frac{(\SI{10e-6}{\metre})^2}{\SI{1.4e-10}{\metre\squared\per\second}} = \SI{0.7}{\second}.
\}
Dit is een typische tijdschaal for oxygen- en nutri\"enttransport in weefsels.
\end{voorbeeld}

\section{Diffusie-gelimiteerde reacties}

In Hoofdstuk~7 behandelden we kinetiek vanuit de snelheidswet $r = k[A]$. Voor snelle reacties kan echter de diffusie van reactanten naar elkaar (of naar een katalysatoroppervlak) de snelheidsbepalende stap worden. De \emph{diffusie-gelimiteerde} snelheidsconstante voor een bimoleculaire reactie in oplossing volgt uit de theorie van Smoluchowski:
\begin{equation}
k_D = 4\pi N_A (D_A + D_B) (r_A + r_B).
\end{equation}
Voor typische kleine moleculen in water ($D \sim \SI{1e-9}{\metre\squared\per\second}$, $r \sim \SI{0.5}{\nano\metre}$) bedraagt $k_D \sim \SI{6e9}{\litre\per\mole\per\second}. Reacties met snelheidsconstanten in deze orde worden volledig door diffusie bepaald: every encounter leidt tot reactie.

\begin{voorbeeld{Snelheid van enzym--substraat-ontmoeting}
In water is de diffusie-gelimiteerde snelheidsconstante ongeveer \SI{1e9}{\litre\per\mole\per\second}. Een enzym met $k_{\text{cat}}/K_M$ van \SI{1e8}{\litre\per\mole\per\second} is reeds dicht bij deze limiet --- het enzym is \emph{diffusie-perfect}: substraat dat het enzym bereikt, wordt vrijwel zeker omgezet. Voorbeelden: carbonanhydrase, triosefosfaatisomerase.
\end{voorbeeld}

\section{Grenzelaagconcepten}

Wanneer een vloeistof langs een vaste wand stroomt, is de snelheid nul aan de wand (no-slip-conditie). Er ontstaat een \emph{snelheidsgrenslaag} waarin de snelheid toenimt van nul aan de wand tot de bulksnelheid. Analoog ontstaan concentratie- en temperatuursgrenslagen wanneer er massaoft warmteoverdracht is tussen wand en bulk.

De grenslaagdikte $\delta$ hangt af van het stromingsregime (laminaal of turbulent) en de transportco\"effici\"enten. Voor laminaire stroming is $\delta \propto \sqrt{\nu x / v_\infty}$, waarin $\nu = \eta/\rho$ de kinematische viscositeit en $x$ de afstand langs de wand is.

De \emph{massaoverdrachtsco\"effici\"ent} $k_L$ (in \si{\metre\per\second}) beschrijft de flux per concentratieverschil:
\begin{equation}
j = k_L\,(c_{\text{bulk}} - c_{\text{wand}}).
\end{equation}
Voor een grenslaag van dikte $\delta$ geldt $k_L \approx D/\delta$. Praktische correlaties ( bv. van Sherwood, Reynolds, Schmidt) worden gebruikt in equipmentontwerp.

\section{Elektrische transportverschijnselen: de wet van Ohm als thermodynamisch verschijnsel}

Elektrische ladingtransport is een bijzonder geval van transportverschijnselen. De wet van Ohm $j = -\sigma_e\,\nabla\phi$ heeft dezelfde vorm als Fick, Fourier en Newton, met elektrische conductiviteit $\sigma_e$ en potentiaalgradi\"ent $\nabla\phi$. De entropyproductie per eenheidsvolume is:
\begin{equation}
\sigma_{\text{entropie}} = \frac{j^2}{\sigma_e T} \geq 0.
\end{equation}
Dit verklaart whyheerweerstand in een elektrische draad warmte produceert (Joule-verwarming): de entropyproductie is altijd positief.

\subsection{De Einstein-relatie}

Voor geladen deeltjes in dilute oplossing is er een fundamenteel verband tussen mobiliteit $\mu_e$ (snelheid per eenheid elektrisch veld) en diffusieco\"effici\"ent:
\begin{equation}
\boxed{\;D = \mu_e\,\frac{k_B T}{e}\;}
\end{equation}
Dit is de \emph{Einstein-relatie}, die aantoont dat de spontane diffusie en de door het elektrisch veld gedwongen beweging dezelfde moleculaire friction voelen. Een rechtstreekse verbinding met de wetten van Nernst en Ohm.

\section{De brug met thermodynamica}

In Hoofdstuk~3 zagen we dat $\Delta G < 0$ voor spontane processen bij constante $T, p$. De flux van materie of warmte is ook een spontaan proces, gedreven door een gradi\"ent in chemische potentiaal of temperatuur. De entropyproductie per volume in een systeem met concentratiegradi\"ent kan worden geschreven als:
\begin{equation}
\sigma_{\text{entropie}} = -J \cdot \nabla\mu / T,
\end{equation}
waarin $\nabla\mu$ de gradi\"ent in chemische potentiaal is. Voor ideale mengsels geldt $\nabla\mu = RT\,\nabla\ln c$, zodat met $J = -D\nabla c$ direct Fick volgt --- diffusie is dus een direct gevolg van het streven naar minimal Gibbs-energie.

\begin{besluit}{Transport: de snelheid tussen twee evenwichten}
Transportverschijnselen vormen de schakel tussen thermodynamica en procestechnologie. De thermodynamica bepaalt wat de evenwichtstoestand is; transportverschijnselen bepalen hoe snel die toestand wordt bereikt. De drie wetten (Fourier, Fick, Newton) delen een gemeenschappelijke wiskundige structuur en een gemeenschappelijke thermodynamische oorsprong (entropyproductie). Voor de ingenieur is de karakteristieke tijdschaal $t \sim L^2/D$ een onmisbaar instrument voor het schatten van procesdynamiek. Diffusie-gelimiteerde reacties vormen een directe brug met Hoofdstuk~7, terwijl de entropyproductie een directe brug vormt met Hoofdstuk~3.
\end{besluit}
